\section{A plan's capacity is different from an employee's award}
\label{sec:apple-stock-plan}

The share register has already separated awards, vesting and common shares
issued. The documents governing an award explain why those quantities can move
differently. Apple's 4 March 2022 Form 8-K supplies the 2022 Employee Stock Plan,
a time-based restricted-stock-unit form and a performance-award form as
Exhibits 10.1--10.3. These are the counterparts of three references in the
annual report. The filing reports shareholder approval and effectiveness on
4 March 2022. The plan also records Board adoption on 9 November 2021 and uses
that earlier date as its defined Effective Date when setting its term. The two
dates perform different jobs; replacing them with one undifferentiated date
would lose information needed to interpret the plan
\citep[Items 5.02 and 5.07; Exhibit 10.1, Sections 1 and 15(a)]{apple2022plan}.

Start with the capacity to make awards. Section 5(a) supplies 510 million
shares plus specified amounts available or recycled from the 2014 plan, subject
to a maximum of 1,274,374,682. That maximum is neither the current unused
capacity nor a promise to issue that many common shares. Section 5(b) counts
a full-value award at two reserve shares for each underlying share, whereas
exercised options and stock appreciation rights use one reserve share for each
share to which the exercise relates. Options and stock appreciation rights
count their gross underlying amount even when fewer shares are delivered
\citep[Exhibit 10.1, Section 5(a)--(b)]{apple2022plan}.

A small hypothetical award makes the distinction concrete. Suppose 100
restricted stock units were initially charged against the reserve at two for
one, all 100 vest and 20 shares are withheld for employee taxes. Assume no
other adjustments. The employee receives 80 shares. Under the plan's
full-value withholding rule, the 20 withheld shares return 40 reserve shares,
because the original charge used that ratio. The two quantities are therefore
\[
 \underbrace{100-20=80}_{\text{common shares delivered}},\qquad
 \underbrace{2(100)-2(20)=160}_{\text{net reserve shares used}}.
\]
The reserve records authority used, rather than shares in the employee's
account. An option exercise covering 100 shares can also deliver 80 after
withholding, yet uses 100 reserve shares under the gross-count rule; the
option withholding does not restore capacity. The same observed delivery
therefore cannot identify the reserve movement without the award type and
counting rule. Cash settlement creates another distinction: the plan excludes
the shares that would otherwise have been delivered from its share count.
These are consequences of Section 5(b), not an assumed employee tax rate,
cash payment or compensation expense \citep[Exhibit 10.1, Section 5(b)]{apple2022plan}.

The 2025 proxy illustrates why the reporting population must travel with each
number. At 28 September 2024 its equity-compensation table reports 125,973,077
shares subject to outstanding awards under the 2022 plan, 37,387,521 under the 2014 plan and
10,612 under the director plan, totalling 163,371,210. The 2014 amount includes
39,984 options and 37,347,537 restricted stock units. Removing those options
leaves 163,331,226 shares subject to restricted stock units across the three listed populations.
The annual note's 163,326 thousand restricted stock units is a different
reported amount; these disclosures alone do not identify a reconciliation.
The proxy separately reports 1,022,428,528 shares available under the 2022
plan, 66,924,428 under the employee purchase plan and 4,110,122 under the
director plan, totalling 1,093,463,078. That total includes different plans
from the outstanding-award total and is governed by their counting conventions
\citep[printed p.~97 and table footnotes]{apple2025proxy}. A reconstruction
must retain these populations rather than subtracting unlike totals to infer
grants, dilution or unused authority.

Capacity still does not tell the employee what will be delivered or when.
Section 9 permits restricted-stock-unit agreements to differ: the agreement
specifies vesting, settlement may be in shares, cash or a combination, and
delivery may be deferred under its terms. Dividend-equivalent rights may be
granted; the units themselves carry no voting rights. The filed time-based
form then makes a more specific promise. Its Section 7 provides for shares
equal to vested units less the applicable tax-related withholding, delivered
promptly and no later than two and one-half months after vesting, subject to
the stated conditions. The plan's wider permission does not replace that
form's narrower promise \citep[Exhibit 10.1, Section 9; Exhibit 10.2, Sections
7 and 11]{apple2022plan}.

The form's grant notice leaves the recipient, quantity, award date and vesting
schedule blank. Those blanks matter: the annual report's general four-year
practice cannot fill every individual notice. For example, the 2025 proxy
describes named executives' 1 October 2023 time-based awards as vesting in
three equal annual instalments beginning 1 April 2026, over approximately
four and one-half years. It also describes termination terms under which
earned or vested rights and delivery dates must be kept separate: Cook's
retirement treatment does not accelerate the originally scheduled delivery
of shares or dividend equivalents
\citep[Exhibit 10.2, grant notice and Section 8]{apple2022plan}
\citep[printed pp.~51 and 53--54]{apple2025proxy}. These examples establish
that a general vesting convention cannot substitute for the applicable award;
they do not identify every employee's terms or a complete amendment history.

For the simulation, the consequence is concrete. Retain the plan version and
reserve rule alongside the actual award terms, vesting conditions, units
earned, settlement medium, delivery date and withholding choice. A permitted
plan alternative becomes a branch for a particular award only when its terms
or explicitly stated unresolved facts allow it. Otherwise cash-settlement
paths would be invented for a share-only promise, or vesting would create
shares before their contractual delivery. Tax amounts, accounting expense
and diluted EPS require their own facts and methods. The performance-award
form poses the next question: how are units earned when the contractual
performance schedule itself has not been filled in?
